% Two state machines side by side: the transaction shape and the repaired shape.
\begin{tikzpicture}[font=\sffamily\scriptsize]

% ================================================= left: the transaction shape
\node[font=\sffamily\bfseries\small, anchor=south] at (1.9,3.3)
  {The transaction shape};
\node[initial] (i1) at (0,2.2) {};
\node[state, text width=20mm] (ready) at (1.9,2.2) {Ready\\{\tiny receive interrupt off}};
\node[state, text width=20mm] (recv) at (1.9,-0.4) {Receiving\\{\tiny counting down n}};
\draw[trans] (i1) -- (ready);
\draw[trans] (ready.west) to[out=200,in=160,looseness=1.2]
  node[left, align=right, xshift=-1mm]{start\\receive(n)} (recv.west);
\draw[trans] (recv.east) to[out=20,in=-20,looseness=1.2]
  node[right, align=left, xshift=1mm]{n reaches 0\\interrupt off} (ready.east);

\node[umlnote, text width=38mm, anchor=north] at (1.9,-1.9)
  {Failure mode 1. Every byte that arrives while the machine sits in Ready is
   lost, and nothing counts it. The window is a fixed length of time, so the
   fraction lost grows with the line rate.};

% ================================================= right: the repaired shape
\node[font=\sffamily\bfseries\small, anchor=south] at (9.5,3.3)
  {The repaired shape};
\node[initial] (i2) at (7.4,2.2) {};
\node[state, text width=20mm] (en) at (9.5,2.2) {Enabled\\{\tiny and never disabled}};
\node[state, text width=20mm] (ovr) at (9.5,-0.4) {Overrun seen};
\draw[trans] (i2) -- (en);
\draw[trans] (en.north east) to[out=40,in=-40,looseness=5] (en.south east);
\node[anchor=west, align=left, font=\sffamily\tiny] at (11.6,2.2)
  {RXNE: drain the\\whole FIFO into\\the ring buffer};
\draw[trans] (en.west) to[out=200,in=160,looseness=1.2]
  node[left, align=right, xshift=-1mm]{ORE sets} (ovr.west);
\draw[trans] (ovr.east) to[out=20,in=-20,looseness=1.2]
  node[right, align=left, xshift=1mm]{write ORECF\\ovr++} (en.east);
\draw[irq] (ovr.south west) to[out=250,in=290,looseness=3.5] (ovr.south east);
\node[anchor=north, align=center, font=\sffamily\tiny, text=red!60!black]
  at (9.5,-1.85) {failure mode 2: return without clearing,\\and this loop never ends};

\node[umlnote, text width=38mm, anchor=north] at (9.5,-3.1)
  {The repaired machine has one resting state. There is no count, no completion
   and no re-arming, so there is no window. An overrun is a counter that goes
   up and not an error that ends anything.};
\end{tikzpicture}
